\begin{tabular}{l|l|l}
\textit{Begrepp} & \textit{Engelska} & \textit{Beskrivning} \\ \hline \hline
@main & @main & där exekveringen av kompilerat program startar \\
abstrahera & abstract & att införa nya begrepp som förenklar kodningen \\
abstrakt klass & abstract class & kan ha parametrar, kan ej instansieras, kan ej mixas in \\
abstrakt medlem & abstract member & saknar implementation \\
algoritm & algorithm & stegvis beskrivning av en lösning på ett problem \\
anonym funktion & anonymous function & funktion utan namn; kallas även lambda \\
anonym klass & anonymous class & klass utan namn, utvidgad med extra implementation \\
Array & Array & en förändringsbar, indexerbar sekvenssamling \\
arv & inheritance & arv beskriver relationen 'är en' \\
attribut & attribute & variabel som utgör (del av) ett objekts tillstånd \\
bastyp & base type & den mest generella typen i en arvshierarki \\
block & block & kan ha lokala namn; sista raden ger värdet \\
boolesk & boolean & antingen sann eller falsk \\
case-klass & case class & slipper skriva new; automatisk innehållslikhet \\
datastruktur & data structure & många olika element i en helhet; elementvis åtkomst \\
dator & computer & en elektronisk enhet som behandlar data enligt instruktioner \\
de-serialisera & deserialize & avkoda symbolsekvens och återskapa objekt i minnet \\
defaultargument & default argument & gör att argument kan utelämnas \\
dynamisk bindning & dynamic binding & körtidstypen avgör vilken metod som körs \\
element & element & objekt i en datastruktur \\
enhetlig access & uniform access & ändring mellan def och val påverkar ej användning \\
enumeration & enumeration & en uppräkning av värden i en viss ordning \\
exekveringsfel & runtime error & kan inträffa medan programmet kör \\
export & export & gör namn synligt utåt som medlem i detta objekt \\
fabriksmetod & factory method & hjälpfunktion för indirekt konstruktion \\
flyttal & floating-point number & decimaltal med begränsad noggrannhet \\
for-sats & for statement & bra då antalet repetitioner är bestämt i förväg \\
funktion & function & vid anrop beräknas ett returvärde \\
funktionshuvud & function header & har parameterlista och eventuellt en returtyp \\
funktionskropp & function body & koden som exekveras vid funktionsanrop \\
förseglad typ & sealed type & subtypning utanför denna kodfil är förhindrad \\
generisk & generic & har abstrakt typparameter, typen är generell \\
getter & getter & indirekt åtkomst av attributvärde \\
implementation & implementation & en specifik realisering av en algoritm \\
import & import & gör namn tillgängligt lokalt utan att hela sökvägen behövs \\
inmixning & mixin & tillföra flera egenskaper genom arv av trait \\
innehållslikhet & structural equality & instanser anses lika om de har samma tillstånd \\
instans & instance & upplaga av ett objekt med eget tillståndsminne \\
klass & class & en mall för att skapa flera instanser av samma typ \\
klassparameter & class parameter & binds till argument som ges vid konstruktion \\
kolonn & column & annat ord för kolumn \\
kolumnvektor & column vector & matris av dimension $m\times{}1$ med $m$ vertikala värden \\
kompanjonsobjekt & companion object & ser privata medlemmar i klass med samma namn \\
kompilator & compiler & ett program som tar källkod som indata och ger maskinkod som utdata \\
kompilera & compile & att översätta kod till exekverbar form \\
kompilering & compilation & kod på lägre nivå skapas ur källkodsfiler \\
kompileringsfel & compile error & kan inträffa innan exekveringen startat \\
konstruktor & constructor & skapar instans, allokerar plats för tillståndsminne \\
kontrollstruktur & control structure & kan ge alternativa vägar genom koden \\
körtidstyp & runtime type & kan vara mer specifik än den statiska typen \\
lat initialisering & lazy initialization & allokering sker först när namnet refereras \\
linjärsöka & linear search & leta i sekvens tills sökkriteriet är uppfyllt \\
linjärsökning & linear search algorithm & sökalgoritm som letar i sekvens tills element hittas \\
litteral & literal & anger ett specifikt datavärde \\
map & map & applicerar en funktion på varje element i en samling \\
mappning & mapping & nyckel -> värde \\
matris & matrix & indexerbar datastruktur i två dimensioner \\
medlem & member & tillhör ett objekt; nås med punktnotation om synlig \\
metod & method & funktion som är medlem av ett objekt \\
minneskomplexitet & memory complexity & hur minnesåtgången växer med problemstorleken \\
modul & module & kodenhet med abstraktioner som kan återanvändas \\
mängd & set & oordnad samling med unika element \\
namnanrop & call by name & fördröjd evaluering av argument \\
namngivna argument & named arguments & gör att argument kan ges i valfri ordning \\
namnrymd & namespace & omgivning där är alla namn är unika \\
namnskuggning & name shadowing & lokalt namn döljer samma namn i omgivande block \\
new & new & nyckelord vid direkt instansiering av klass \\
null & null & ett värde som ej refererar till någon instans \\
nyckel & key & en unik identifierare \\
nyckel-värde-tabell & key-value table & oordnad samling av mappningar med unika nycklar \\
objekt & object & samlar variabler och funktioner \\
ordning & ordering & definierar hur element av en viss typ ska ordnas \\
paket & package & modul som skapar namnrymd; maskinkod får egen katalog \\
parameterlista & parameter list & beskriver namn och typ på parametrar \\
persistens & persistence & egenskapen att finnas kvar efter programmets avslut \\
polymorfism & polymorphism & kan ha många former, t.ex. en av flera subtyper \\
predikat & predicate & en funktion som ger ett booleskt värde \\
privat & private & modifierar synligheten av en objektmedlem \\
procedur & procedure & vid anrop sker (sido)effekt; returvärdet är tomt \\
programargument & program argument & kan överföras via parametern args till main \\
punktnotation & dot notation & används för att komma åt icke-privata delar \\
radvektor & row vector & matris av dimension $1\times{}m$ med $m$ horisontella värden \\
Range & Range & en samling som representerar ett intervall av heltal \\
referenslikhet & reference equality & instanser anses olika även om tillstånden är lika \\
referenstyp & reference type & har supertypen AnyRef, allokeras i heapen via referens \\
registrering & counting & algoritm som räknar element med vissa egenskaper \\
rekursiv funktion & recursive function & en funktion som anropar sig själv \\
samling & collection & datastruktur med element av samma typ \\
samlingsbibliotek & collection library & många färdiga samlingar med olika egenskaper \\
sats & statement & en kodrad som gör något; kan särskiljas med semikolon \\
sekvens(samling) & sequence (collection) & noll el. flera element av samma typ i viss ordning \\
sekvensalgoritm & sequence algorithm & lösning på problem som drar nytta av sekvenssamling \\
sekvenssamling & sequence collection & datastruktur med element i en viss ordning \\
serialisera & serialize & koda objekt till avkodningsbar sekvens av symboler \\
setter & setter & indirekt tilldelning av attributvärde \\
singelobjekt & singleton object & modul som kan ha tillstånd; finns i en enda upplaga \\
skript & script & ensam kodfil, huvudprogram behövs ej \\
skyddad medlem & protected member & är endast synlig i subtyper \\
slumptalsfrö & random seed & ger återupprepningsbar sekvens av pseudoslumptal \\
sortering & sorting & algoritm som ordnar element i en viss ordning \\
stack trace & stack trace & lista anropskedja vid körtidsfel \\
sträng & string & en sekvens av tecken \\
stränginterpolator & string interpolator & en funktion för att bädda in uttryck i strängar \\
subtyp & subtype & en typ som är mer specifik \\
supertyp & supertype & en typ som är mer generell \\
sökning & search & algoritm som letar upp element enligt sökkriterium \\
tidskomplexitet & time complexity & hur exekveringstiden växer med problemstorleken \\
tilldelning & assignment & för att ändra en variabels värde \\
trait & trait & är abstrakt, kan mixas in, kan ha parametrar \\
typ & type & beskriver vad data kan användas till \\
typalias & type alias & alternativt namn på typ som ofta ökar läsbarheten \\
typargument & type argument & konkret typ, binds till typparameter vid kompilering \\
typhärledning & type inference & kompilatorn beräknar typ ur sammanhanget \\
uttryck & expression & kombinerar värden och funktioner till ett nytt värde \\
variabel & variable & kan initialiseras, refereras och eventuellt ändras med tilldelning \\
Vector & Vector & en oföränderlig, indexerbar sekvenssamling \\
värdeanrop & call by value & argumentet evalueras innan anrop \\
värdetyp & value type & har supertypen AnyVal, lagras direkt på stacken \\
while-sats & while statement & bra då antalet repetitioner ej är bestämt i förväg \\
yield & yield & används i for-uttryck för att skapa ny samling \\
äkta funktion & pure function & ger alltid samma resultat om samma argument \\
överlagring & overloading & metoder med samma namn men olika parametertyper \\
överskuggad medlem & overridden member & medlem i subtyp ersätter medlem i supertyp \\
\end{tabular}
